% === I. Judul / Nama Perangkat Lunak ===
\secbox{I.}{Judul / Nama Perangkat Lunak}

\begin{abstractbox}
\centering
\textbf{LUNA (Listening, Understanding, Nurturing Assistant): Platform Pendamping Kesehatan Mental Berbasis Multi-Agent Artificial Intelligence dengan Personalized Mental Health Support, Automatic Diary, dan Early Mental Crisis Detection bagi Generasi Z}
\end{abstractbox}

\indent
\textbf{LUNA (Listening, Understanding, Nurturing Assistant)} merupakan platform pendamping kesehatan mental berbasis \textit{Multi-Agent Artificial Intelligence} yang dirancang untuk memberikan dukungan kesehatan mental secara personal bagi Generasi Z. Nama LUNA sendiri merepresentasikan kemampuan sistem dalam mendengarkan (\textit{listening}), memahami kondisi emosional pengguna (\textit{understanding}), memberikan dukungan yang adaptif (\textit{nurturing}), serta berperan sebagai asisten digital (\textit{assistant}) yang siap mendampingi pengguna melalui percakapan berbasis teks maupun suara.

Selain itu, LUNA mengintegrasikan beberapa agen kecerdasan buatan yang bekerja secara kolaboratif untuk menyediakan \textit{Personalized Mental Health Support}, menyusun \textit{Automatic Diary} berdasarkan hasil percakapan, serta melakukan \textit{Early Mental Crisis Detection} sebagai upaya deteksi dini terhadap potensi krisis kesehatan mental.

\vspace{0.5em}
\begin{abstractbox}
\centering
\textbf{Tautan Akses Berkas, Source Code, \& Video Demo Perangkat Lunak:}\\[0.3em]
\small
\url{https://drive.google.com/drive/folders/1bJp6CdtFBXfh3RJ8tv2wlBQT3z1xdhio?usp=drive_link}\\[0.2em]
{\small\itshape (Folder Google Drive berisi Berkas Source Code, Build APK, Dokumentasi Teknis, dan Video Demo Aplikasi LUNA)}
\end{abstractbox}
